\documentclass[11pt]{article}

% ---- packages (all standard; compiles with pdflatex) ----
\usepackage[utf8]{inputenc}
\usepackage[T1]{fontenc}
\usepackage[margin=1in]{geometry}
\usepackage{amsmath,amssymb}
\usepackage{booktabs}
\usepackage{graphicx}
\usepackage{caption}
\usepackage{microtype}
\usepackage[hidelinks]{hyperref}

\title{\textbf{Multimodal Deep Fusion of Genomic Sequence, DNA Shape,\\
and Biochemical Context for SP-Family Transcription-Factor\\
Binding-Site Prediction: a Leakage-Controlled Benchmark}}
\author{}
\date{}

\begin{document}
\maketitle

\begin{abstract}
\noindent
Mapping transcription-factor (TF) binding sites is central to understanding gene
regulation. The SP family (SP1, SP2, SP4) recognises a nearly identical GC-box,
which makes per-paralog prediction in short windows unreliable. We build a
multimodal architecture, \emph{G-CMAB}, that fuses three views of a 101\,bp window:
semantic sequence embeddings from the pre-trained genomic transformer DNABERT-2,
local three-dimensional geometry from DNAshape, and biochemical descriptors
(CpG observed/expected ratio, GC content, G-quadruplex motifs). The sequence and
shape streams are coupled by bidirectional cross-modal attention and a multi-scale
depthwise-separable CNN with DDP-safe Group Normalization. After an evidence-driven
pivot from four-way to binding-versus-non-binding classification, the model attains
$81$--$82\%$ accuracy (ROC-AUC $0.83$--$0.86$) on a leakage-controlled benchmark
with composition-matched dinucleotide-shuffled negatives, modestly above a
fine-tuned sequence-only DNABERT-2 ($78.7\%$) and a strong classical TF-IDF
logistic-regression baseline ($80.6\%$), and consistent ($\sim\!80\%$) on real
GC-matched genomic negatives. We additionally report a cautionary methodological
result: an earlier $96.35\%$ figure was traced to a DNAshape sequence--shape
pairing artifact and is withdrawn. Threshold calibration raises Negative recall to
$0.57$ at $79.6\%$ accuracy. We release all data-handling, training, ablation, and
interpretability code.
\end{abstract}

\section{Introduction}

\subsection{Biological background}
SP-family TFs regulate proliferation, metabolism and apoptosis; their dysregulation
is linked to cancer and cardiovascular disease, so accurate binding-site maps are
valuable for mechanistic study and therapeutic research. Chromatin
immunoprecipitation sequencing (ChIP-seq) yields experimental binding maps, but
computational predictors are cheaper and scale better.

\subsection{Limitations of prior approaches}
Most machine-learning and recent deep-learning predictors operate on the flat
one-dimensional base sequence (A, C, G, T). Even strong pre-trained encoders such as
DNABERT-2 remain extensions of the 1D view and tend to over-predict binding,
ignoring the physical shape of the double helix and the surrounding biochemical
context that influence binding. A second, under-appreciated issue is evaluation
rigour: impressive accuracies often fail to reproduce under stricter data handling.
Our own work provides a concrete instance --- a data-pairing artifact that inflated
accuracy by roughly $15$ points (Section~\ref{sec:artifact}).

\subsection{Contributions}
\begin{enumerate}
  \item A complete, reproducible preprocessing pipeline: ENCODE ingestion,
        cross-target overlap removal, class balancing, and composition-preserving,
        leakage-free negative construction.
  \item A multimodal tri-branch architecture (G-CMAB) that couples DNABERT-2 sequence
        embeddings with DNAshape geometry through bidirectional cross-modal attention
        and a multi-scale CNN, plus a biochemical branch.
  \item A leakage-controlled benchmark and the diagnosis/correction of a DNAshape
        pairing artifact, with results that are stable across negative types and
        train/test splits.
  \item Model interpretability via SHAP attributions and genome-browser (BED) export.
\end{enumerate}

\section{Related Work}
Sequence-only motif learners such as DeepBind~\cite{deepbind} established the CNN
baseline. Transformer approaches followed: BERT-TFBS~\cite{berttfbs} and
DNABERT-based multi-scale convolution/attention frameworks couple transformer
embeddings with convolutions, and BERT+2D-CNN hybrids~\cite{le} outperform classical
models. DNABERT-2~\cite{dnabert2} provides efficient multi-species genomic
embeddings, and DNAshapeR~\cite{dnashaper} predicts helix geometry from a
pentamer lookup table. A persistent gap is the lack of explicit cross-interaction
between sequence and 3D structure; we address it with bidirectional cross-modal
attention.

\section{Materials and Methods}

\subsection{Data source and curation}
We used ENCODE HepG2 ChIP-seq narrowPeak tracks on the hg38 assembly (SP1
ENCFF333SWC, SP2 ENCFF480YAW, SP4 ENCFF938KVY), chosen for consistent cell line and
assay. Peaks were filtered at $-\log_{10}(q)\ge 2.0$ (i.e.\ $q\le 0.01$). Because
paralogous SP factors bind overlapping loci, we assigned each overlapping region to
its highest-$q$ target, yielding $6{,}815$ / $1{,}696$ / $8{,}914$ exclusive peaks for
SP1 / SP2 / SP4. We down-sampled to the SP2 baseline ($1{,}696$ per TF) by descending
$q$-value to balance the classes.

\subsection{Sequence extraction, augmentation and splitting}
Fixed $101$\,bp windows were centred on each summit ($\pm 50$\,bp); boundary-violating
windows were discarded and bases upper-cased. Positives were doubled by
reverse-complement (RC) augmentation, giving $3{,}392$ per TF and $13{,}568$ sequences
in total. Data were split $80/20$ with a grouped split that keeps each
(original, RC) pair in the same fold, preventing augmentation leakage; a stricter
chromosome-holdout protocol was used for robustness checks.

\subsection{Negative-set construction}
The primary negatives are \emph{dinucleotide-shuffled} sequences generated with the
Altschul--Erickson Eulerian-path algorithm~\cite{altschul}, which preserves mono-
and dinucleotide frequencies (hence GC and CpG density) exactly while destroying
motif structure. This is a deliberately hard, composition-matched control: a model
cannot exploit base composition as a shortcut. For corroboration we also built
GC- and length-matched real genomic negatives (ENCODE blacklist exclusion,
$\pm 500$\,bp peak padding, $2\%$ GC bins, $<\!5\%$ N, Cohen's-$d$ quality control).

\subsection{Feature engineering (three modalities)}
\textbf{(a) Sequence.} Windows are tokenised with the DNABERT-2 byte-pair vocabulary;
the maximum token length is set to the training-set $99$th percentile (clamped to
$[32,96]$). \textbf{(b) DNAshape.} Five structural parameters --- Minor Groove Width,
Propeller Twist, Roll, Helical Twist, Electrostatic Potential --- are computed per
position by a sliding pentamer lookup, with reverse-complement fallback and edge-NaN
handling, producing an $[N,5,101]$ tensor that is robustly scaled (median-centred,
divided by the $1$st--$99$th-percentile range, stats from the training split only).
\textbf{(c) Biochemical.} CpG O/E ratio, GC content, and G-quadruplex presence
(regular-expression match) per window.

\subsection{Architecture: G-CMAB}
The three branches feed a fusion head. The \emph{sequence} branch unfreezes the last
six DNABERT-2 encoder layers and applies an ELMo-style learned scalar-mix over the
last hidden states, projected to $128$ dimensions. The \emph{shape} branch projects
the $[N,5,101]$ tensor to $128$ channels (1$\times$1 convolution, GroupNorm, GELU).
A bidirectional \emph{cross-modal multi-head attention} module ($d_{\text{model}}=128$,
$4$ heads) lets sequence query shape and shape query sequence, aligning nucleotide
context to helix geometry. Both streams then pass a \emph{multi-scale
depthwise-separable CNN} (sequence kernels $\{7,9,11,15\}$, shape kernels
$\{4,8,12,16\}$, $256$ channels, 1-max pooling). The \emph{bio} branch is a small
MLP ($3{\to}16{\to}32$). The pooled sequence, pooled shape, and bio vectors are
concatenated ($2{,}080$-d) and classified by a sigmoid head. The umbrella name
G-CMAB denotes \textbf{G}roupNorm $+$ \textbf{C}ross-\textbf{M}odal
\textbf{A}ttention $+$ multi-scale \textbf{B}ranch.

\subsection{Training and optimisation}
We minimise \texttt{BCEWithLogits} loss with a dynamic positive-class weight to
counter the $3{:}1$ positive:negative ratio, AdamW with branch-specific learning
rates (backbone $2\times 10^{-5}$; projections/attention $2$--$3\times 10^{-4}$;
layer-attention $10^{-3}$), weight decay $0.1$, linear warmup ($\approx\!15\%$) then
cosine annealing, gradient clipping, and bf16 mixed precision. GroupNorm replaces
BatchNorm throughout the convolutional branches because per-GPU batches are small
under distributed data parallelism. Critically, the data loaders \emph{enforce}
that the negative sequence file and the negative DNAshape file come from the same
set, refusing any mismatched fallback (Section~\ref{sec:artifact}). At evaluation we
report metrics at the default $0.5$ threshold and, additionally, scan the decision
threshold to expose the precision/recall operating curve.

\section{Results}

\subsection{The multiclass plateau and the binary pivot}
Every four-way model (SP1 / SP2 / SP4 / Negative) saturated between $\sim\!50$ and
$62\%$ accuracy. The error pattern was consistent: the Negative class was recovered
well, but the three SP paralogs were mutually confused (SP4 recall as low as
$0.33$--$0.57$). A $101$\,bp window therefore answers ``is this a binding site?'' but
not ``which SP paralog?'', because the paralogs share motif and shape signature. We
read this as a property of the data and reformulated the task as binding
(SP1$\cup$SP2$\cup$SP4) versus non-binding.

\subsection{Binary classification}
Table~\ref{tab:binary} reports the binary comparison on the clean, sequence--shape
matched dinucleotide benchmark. The multimodal G-CMAB reaches $81.0$--$82.3\%$
accuracy across runs and configurations (ROC-AUC $0.83$--$0.86$), modestly above the
fine-tuned sequence-only DNABERT-2 ($78.7\%$) and the best classical model
(TF-IDF logistic regression, $80.6\%$). The improvement over sequence-only is real
but small and shows up most clearly in ranking (ROC-AUC) rather than in
default-threshold accuracy.

\begin{table}[ht]
\centering
\caption{Binary classification on the clean dinucleotide benchmark (held-out test,
$n=2{,}710$). Deep rows are verified sequence--shape-matched runs; the multimodal
row gives the instrumented reference run (range $81.0$--$82.3\%$ across
configurations). Classical baselines ($\dagger$) default to the genomic GC-matched
negatives and are within $\sim\!0.2$ points on dinucleotide negatives.}
\label{tab:binary}
\begin{tabular}{lccccc}
\toprule
Model & Acc. & ROC-AUC & PR-AUC & Neg.\ R & Pos.\ R \\
\midrule
Classical one-hot $\cdot$ Gradient Boosting$^{\dagger}$ & 0.755 & 0.624 & 0.823 & 0.043 & 0.993 \\
Sequence-only DNABERT-2 & 0.787 & 0.817 & 0.926 & 0.585 & 0.855 \\
Classical TF-IDF $\cdot$ Logistic Regression$^{\dagger}$ & 0.806 & 0.818 & 0.928 & 0.376 & 0.951 \\
\textbf{Multimodal G-CMAB} & \textbf{0.810} & \textbf{0.832} & \textbf{0.919} & 0.402 & 0.937 \\
\quad G-CMAB, enhanced (2 attn.\ layers, 8 layers unfrozen) & 0.823 & --- & --- & --- & --- \\
\quad G-CMAB @ tuned threshold ($t=0.95$) & 0.796 & 0.832 & 0.919 & \textbf{0.569} & 0.871 \\
\bottomrule
\end{tabular}
\end{table}

\subsection{Threshold calibration}
At the default $0.5$ threshold the multimodal model over-predicts the positive class,
and its default-threshold Negative recall varies run-to-run ($0.40$--$0.61$) while
ROC-AUC is stable ($\sim\!0.83$--$0.86$). Because the ranking is informative, simple
threshold calibration recovers minority-class recall: raising the threshold to
$t=0.95$ lifts Negative recall to $0.569$ while overall accuracy stays at $79.6\%$,
indicating the model has learned a genuine signal rather than guessing the majority
class.

\subsection{Consistency across negative types and splits}
The result is not specific to dinucleotide negatives. Re-evaluated against real
GC- and length-matched genomic negatives, the same architecture scores a consistent
$79.6\%$ (random group split) and $79.4\% \pm 0.8\%$ under a chromosome-holdout,
multi-seed protocol (test chromosomes chr1/6/7; seeds 42/1/2023), so the figure is
stable to both negative type and to genomic-proximity/homology leakage. The
dinucleotide and genomic settings agree once the data pairing is correct.

\subsection{A DNAshape pairing artifact (cautionary result)}
\label{sec:artifact}
An earlier version of this model reported $96.35\%$ accuracy. We traced this to a
data-handling bug: the negative \emph{sequences} were the dinucleotide shuffles, but
the negative \emph{DNAshape} tensor was loaded (by an independent recursive search)
from the \emph{genomic} negative set. Because genomic and shuffled sequences have
markedly different shape profiles, the shape branch could separate the two classes
on a feature that does not correspond to the sequences the rest of the model saw ---
an artificial leak. After enforcing sequence--shape pairing and re-running cleanly,
accuracy falls to $\sim\!81\%$, matching the genomic and classical results. The
project's own corrected real-negative evaluation records the same $\sim\!18$-point
correction. We report this prominently because composition-matched and
geometry-derived negatives are increasingly common in TFBS prediction, and silent
cross-modal mismatches are easy to introduce and hard to notice.

\subsection{Ablation and an honest note on training fragility}
We ran a controlled ablation that toggles individual components (sequence only;
$+$ shape; $+$ cross-attention; $+$ bio; and single-component knockouts) under a
fixed split and seed at $30$ epochs (Table~\ref{tab:abl}). On the clean dinucleotide
data \emph{all} configurable variants --- including the nominally complete one ---
collapsed to the majority class ($74.83\%$ accuracy, the positive fraction;
Negative recall $0.00$; ROC-AUC $\approx 0.52$), whereas the main training pipeline
reached $\sim\!81\%$. We therefore do \emph{not} claim the ablation isolates the
contribution of any single component: a nominally identical ``full'' configuration
collapsed in the ablation harness while the main trainer did not, which points to
sensitivity to the optimisation schedule and training procedure rather than to a
clean architectural effect. The honest reading is that this composition-matched task
sits in a difficult optimisation basin from which under-tuned configurations fail to
escape; isolating component value requires a more carefully tuned ablation, which we
leave to future work.

\begin{table}[ht]
\centering
\caption{Controlled ablation on the clean dinucleotide data ($30$ epochs,
$n=2{,}710$). All configurable variants collapse to the majority class; the table is
reported for transparency, not as clean component attribution (see text).}
\label{tab:abl}
\begin{tabular}{llccc}
\toprule
Variant & Configuration & Acc. & ROC-AUC & Neg.\ R \\
\midrule
A1 & Sequence only            & 0.7483 & $\approx$0.52 & 0.00 \\
A2 & $+$ DNAshape              & 0.7483 & $\approx$0.52 & 0.00 \\
A3 & $+$ Cross-modal attention & 0.7483 & $\approx$0.55 & 0.00 \\
A4 & $+$ Bio (full)           & 0.7483 & $\approx$0.53 & 0.00 \\
A5--A7 & full $-$ one component & $\approx$0.747 & $\approx$0.52 & $\approx$0.00 \\
\midrule
\textbf{Main trainer (script 28)} & full G-CMAB, tuned schedule & \textbf{0.810} & \textbf{0.832} & 0.402 \\
\bottomrule
\end{tabular}
\end{table}

\subsection{Interpretability and genomic output}
SHAP attributions decompose each prediction into sequence, DNAshape and biochemical
contributions around the GC-box. Attributions must be computed on near-decision-boundary
samples: on confidently-correct examples the saturated sigmoid yields near-zero input
gradients, so a gradient-based explainer returns spuriously flat attributions for all
modalities. Correct and confused predictions are exported as BED tracks for
genome-browser inspection and downstream motif/chromatin analysis.

\section{Discussion}

\subsection{What the results say}
On a leakage-controlled benchmark the multimodal model reaches $\sim\!81$--$82\%$,
modestly above a sequence-only transformer and roughly level with a classical TF-IDF
baseline, and consistent across dinucleotide and genomic negatives. The multimodal
edge is real (better ROC-AUC, recoverable Negative recall via thresholding) but
small enough that the architecture is not yet clearly worth its complexity. The most
transferable finding is methodological: a silent sequence--shape pairing mismatch
inflated accuracy by $\sim\!15$ points, a reminder that multimodal genomic pipelines
need explicit modality-alignment checks.

\subsection{Limitations}
(i) Data come from a single cell line (HepG2) and one TF family, so generalisation is
untested. (ii) Absolute performance is modest and only a few points above the
all-positive baseline; lifting it via operating-point calibration and branch upgrades
is ongoing work. (iii) The controlled ablation collapsed to the majority class and is
inconclusive for component attribution. (iv) Per-paralog identification remains
unsolved at $101$\,bp resolution.

\subsection{Performance ceiling and future work}
Accuracy plateaus near $81\%$ across settings, consistent with a $101$\,bp window
holding only so much information once sequence, shape and biochemistry are all used.
Future work includes longer genomic context, chromatin-accessibility signal (e.g.\
ATAC-seq) and quantitative binding affinities in place of binary peak calls, together
with deeper cross-modal fusion and a re-tuned ablation.

\section{Conclusion}
We presented G-CMAB, a tri-branch model that fuses DNABERT-2 sequence embeddings,
DNAshape geometry and biochemical features through bidirectional cross-modal
attention and a multi-scale CNN. After an evidence-driven pivot to binding-versus-
non-binding classification, the model reaches $\sim\!81$--$82\%$ on a leakage-controlled
benchmark, modestly above a sequence-only transformer and stable across negative
types. Equally important, we diagnose and correct a DNAshape pairing artifact that had
inflated accuracy to $96\%$, and we release the full pipeline so the corrected,
honest benchmark is reproducible.

\begin{thebibliography}{9}
\bibitem{dnabert2} Z.\ Zhou \emph{et al.}, ``DNABERT-2: Efficient Foundation Model
and Benchmark for Multi-Species Genome,'' \emph{arXiv:2306.15006}, 2023 (ICLR 2024).
\bibitem{dnashaper} T.-P.\ Chiu \emph{et al.}, ``DNAshapeR: an R/Bioconductor package
for DNA shape prediction and feature encoding,'' \emph{Bioinformatics}
32(8):1211--1213, 2016.
\bibitem{deepbind} B.\ Alipanahi, A.\ Delong, M.\ T.\ Weirauch, B.\ J.\ Frey,
``Predicting the sequence specificities of DNA- and RNA-binding proteins by deep
learning,'' \emph{Nature Biotechnology} 33(8):831--838, 2015.
\bibitem{altschul} S.\ F.\ Altschul, B.\ W.\ Erickson, ``Significance of nucleotide
sequence alignments: a method for random sequence permutation that preserves
dinucleotide and codon usage,'' \emph{Mol.\ Biol.\ Evol.} 2(6):526--538, 1985.
\bibitem{berttfbs} K.\ Wang \emph{et al.}, ``BERT-TFBS: a novel BERT-based model for
predicting transcription factor binding sites by transfer learning,''
\emph{Briefings in Bioinformatics}, 2024.
\bibitem{le} N.\ Q.\ K.\ Le \emph{et al.}, ``A transformer architecture based on BERT
and a 2-D CNN to identify DNA enhancers,'' \emph{Briefings in Bioinformatics}.
\bibitem{shap} S.\ M.\ Lundberg, S.-I.\ Lee, ``A unified approach to interpreting
model predictions,'' \emph{NeurIPS} 30:4765--4774, 2017.
\bibitem{adamw} I.\ Loshchilov, F.\ Hutter, ``Decoupled weight decay
regularization,'' \emph{ICLR}, 2019.
\bibitem{encode} The ENCODE Project Consortium, ``An integrated encyclopedia of DNA
elements in the human genome,'' \emph{Nature} 489(7414):57--74, 2012.
\end{thebibliography}

\end{document}
